\documentclass[12pt,a4paper]{article}
\usepackage[utf8]{inputenc}
\usepackage[brazil]{babel}
\usepackage[top=2.5cm,bottom=2.5cm,left=3cm,right=2.5cm]{geometry}
\usepackage{amsmath,amssymb}
\usepackage{listings}
\usepackage{xcolor}
\usepackage{hyperref}

\hypersetup{
    colorlinks=true,
    linkcolor=black,
    urlcolor=blue,
    citecolor=black
}

\definecolor{codegreen}{rgb}{0,0.6,0}
\definecolor{codegray}{rgb}{0.5,0.5,0.5}
\definecolor{codepurple}{rgb}{0.58,0,0.82}
\definecolor{backcolour}{rgb}{0.96,0.96,0.96}

\lstdefinestyle{mystyle}{
    backgroundcolor=\color{backcolour},   
    commentstyle=\color{codegreen},
    keywordstyle=\color{blue},
    numberstyle=\tiny\color{codegray},
    stringstyle=\color{codepurple},
    basicstyle=\ttfamily\footnotesize,
    breakatwhitespace=false,         
    breaklines=true,                 
    captionpos=b,                    
    keepspaces=true,                 
    numbers=left,                    
    numbersep=5pt,                  
    showspaces=false,                
    showstringspaces=false,
    showtabs=false,                  
    tabsize=4
}

\lstset{style=mystyle}

\begin{document}

% --- FOLHA DE ROSTO ---
\begin{titlepage}
    \centering
    \vspace*{1cm}
    {\Large \textbf{UNIVERSIDADE FEDERAL DO RIO DE JANEIRO}}\\[0.3cm]
    {\large DEPARTAMENTO DE CIÊNCIA DA COMPUTAÇÃO}\\[2.5cm]
    
    {\Large \textbf{RELATÓRIO DE IMPLEMENTAÇÃO}}\\[0.5cm]
    {\LARGE \textbf{E.D.: CLASSE PILHA}}\\[2.5cm]
    
    \begin{flushright}
    \begin{minipage}{0.68\textwidth}
        \textbf{Integrantes e DREs:}\\[0.4cm]
        Gabriel de Souza Machado González -- DRE: 120066108\\[0.2cm]
        Mateus Fernandes Santos -- DRE: 123224894\\[0.2cm]
        Ramon Ramirez -- DRE: \underline{\hspace{3.5cm}}\\[0.2cm]
        Zico Costa -- DRE: 124011595\\[0.8cm]
        \textbf{Repositório Git:}\\
        \url{https://github.com/mateus-ma/classe-pilha.git}
    \end{minipage}
    \end{flushright}
    
    \vfill
    {\large Rio de Janeiro}\\[0.2cm]
    {\large 2026}
\end{titlepage}

\newpage
\pagenumbering{arabic}
% --- CORPO DO RELATÓRIO ---

\section{Análise da Solução Gerada por Inteligência Artificial}

A IA resolveu bem os problemas, apesar de ter apresentado soluções excessivamente técnicas para o contexto do trabalho. Após estudar o código, não encontramos pontos que necessitassem de refatoração, o que faz sentido, pois o problema é bem simples.

Não houve erros de execução, mas o código gerado inicialmente não condizia com a estrutura de repositório sugerida pela própria IA para a implementação em C++, que contava com um arquivo a mais. Um prompt extra resolveu essa questão.

\section{Análise de Trecho Específico em Python}

Um trecho da implementação em Python que chamou atenção foi:

\begin{lstlisting}[language=Python]
    def empilha(self, dado):
        if self.pilha_esta_cheia():
            raise PilhaCheiaErro("A pilha esta cheia. Nao e possivel empilhar mais elementos.")
        
        try:
            self._dados.append(dado)
        except (TypeError, ValueError) as err:
            raise TipoErro(f"O dado '{dado}' nao e compativel com o tipo estipulado '{self._typecode}'.") from err
\end{lstlisting}

A estrutura \texttt{array.array} já exige que todos os dados sejam do mesmo tipo e levanta um \texttt{ValueError} ou \texttt{TypeError} caso um tipo diferente seja passado. Como nas especificações da tarefa foi exigido que o erro exibido seja o customizado \texttt{TipoErro}, a IA foi capaz de perceber e contornar isso.

\section{Resultados dos Testes}

Os testes também foram gerados por IA. Todos os métodos de cada classe foram testados, é possível verificar isso pelo ``OK!'' em C++/JS e os 7 pontos em Python. Repare bem a diferença entre os tempos de execução, aqui, C++ e JavaScript apresentam performances bem semelhantes.

\subsection*{C++}
\begin{verbatim}
Executando testes unitarios (C++)... OK!

[C++ Estresse - 10000000 elementos]
Tempo Empilhar: 0.0392435s
Tempo Desempilhar: 0.0179506s
\end{verbatim}

\subsection*{JavaScript}
\begin{verbatim}
Executando testes unitários (JavaScript)...
OK!

[JS Estresse - 10000000 elementos]
Tempo Empilhar: 0.0385s
Tempo Desempilhar: 0.0190s
\end{verbatim}

\subsection*{Python}
\begin{verbatim}
Tempo de Empilhar: 0.1707s
Tempo de Desempilhar: 0.1288s
.......

Ran 7 tests in 0.300s

OK
\end{verbatim}

\subsection*{Lista de testes em Python:}
\begin{itemize}
    \item Operações básicas;
    \item Teste de exceção para pilha cheia;
    \item Teste de exceção para pilha vazia;
    \item Teste de exceção para tipo incompatível;
    \item Teste do método troca;
    \item Teste de troca quando não há elementos suficientes;
    \item Teste de estresse com 1.000.000 de elementos.
\end{itemize}


\section{Conclusão}

Em suma, a implementação da estrutura de dados Pilha em C++, JavaScript e Python cumpriu com sucesso todos os requisitos propostos. A análise do código gerado pela IA permitiu compreender nuances importantes de implementação: como o tratamento do \texttt{array.array} em Python via bloco \texttt{try-except}. Além disso, pudemos validar a eficácia dos testes unitários, que cobriram todas as nossas demandas. Por fim, os testes de estresse evidenciaram com clareza o comportamento de cada ambiente de execução, destacando a equivalência e superioridade de performance do C++ e do JavaScript frente ao Python para volumes massivos de operações. 
\end{document}